% =====================================================================
% FJCAI 2027 -- paper skeleton.
%
% Built from FJCAI2027_Template_LaTeX.tex. ONLY the four marked blocks
% differ from the shipped template; everything else is byte-identical and
% must stay that way (the template says margins, columns, line spacing and
% fonts are prescribed).
%
% Compile with pdfLaTeX. XeLaTeX/LuaLaTeX WILL FAIL: \usepackage[utf8]{vietnam}
% is a pdfTeX-era package.
%
% Upload beside this file: IEEEtran.cls, fig1_architecture.pdf, fig2_pareto.pdf
% =====================================================================

\documentclass[9pt, conference, a4paper, compsocconf]{IEEEtran}
\usepackage{amsmath,amsxtra,amssymb,amsthm,latexsym,amscd,amsfonts}
\usepackage[utf8]{vietnam}
\usepackage[top=2.7cm, bottom = 5.4cm, left=2.72cm, right = 2.85cm]{geometry}
\headsep = 0.7cm
\headheight = 1.0cm
\footskip = 1.0cm
\voffset = -0.74 cm
%%%%
\usepackage{balance}

\usepackage{fancyhdr}
\pagestyle{fancy}
\usepackage{cite}
\usepackage{algorithmic}
\usepackage{graphicx}
\usepackage{textcomp}
\usepackage{booktabs}
\usepackage{tabularx}
\usepackage{xcolor}
\usepackage{amsmath,amsxtra,amssymb,amsthm,latexsym,amscd,amsfonts}
\usepackage{pgfplots}
\usepackage{multirow}
\usepackage{subcaption}
\pgfplotsset{compat=1.17}

\usepackage{tikz}
\usetikzlibrary{positioning}
\usetikzlibrary{arrows.meta}
\usetikzlibrary{shapes}
\usetikzlibrary{calc}
\usetikzlibrary{fit}
\usepackage[english]{babel}
\usepackage{fancyhdr}
\pagestyle{fancy}
\renewcommand{\sectionmark}[1]{\markright{\MakeUppercase{#1}}{}}
\renewcommand{\headrulewidth}{0pt}
\ifCLASSINFOpdf
\usepackage[pdftex]{graphicx}
  \DeclareGraphicsExtensions{.pdf,.jpeg,.png,.jpg}
\else
\usepackage[dvips]{graphicx}
  \DeclareGraphicsExtensions{.eps}
\fi
\usepackage{array}

\usepackage{caption}
\captionsetup{
  font={footnotesize},
  justification=centering,
  singlelinecheck=false
}
\usepackage[font=footnotesize]{subfig}

% ---------------------------------------------------------------------
% CHANGE 1 of 4: conference year 2026 -> 2027 (appears 3x in this file).
% CONFIRM the day/month and city with the organizers before submitting:
% the shipped template says "Cần Thơ, 27-28/3/2026" and only the YEAR is
% changed below. If FJCAI 2027 moves city or dates, edit those too.
% ---------------------------------------------------------------------
\makeatletter
\def\ps@IEEEtitlepagestyle{%
\def\@oddhead{\hfil \small{\textit{Hội thảo khoa học Quốc gia về Trí tuệ nhân tạo 2027 (FJCAI) - Cần Thơ, 27-28/3/2027}\hfil}%
	\def\@evenhead{\hfil\small{\textit{Hội thảo khoa học quốc gia về Trí tuệ nhân tạo 2027 (FJCAI) - Cần Thơ, 27-28/3/2027}\hfil}}%
		\def\@oddfoot{\scriptsize \thepage \hfil }%
		\def\@evenfoot{\scriptsize \hfil \thepage}
	}
}

\def\@maketitle{%
  \newpage
  \null
  \begin{center}%
    {%
      \fontsize{16}{18}\selectfont
      \bfseries \@title \par
    }%
    \vskip 4em%
    {%
    \fontsize{9}{3.5}\selectfont
      \begin{tabular}[t]{c}%
        \@author
      \end{tabular}\par
    }%
  \end{center}%
  \par
}

\renewenvironment{abstract}
  {\normalfont
   \@IEEEabskeysecsize
   \setlength{\parindent}{0pt}
   \noindent
   \if@twocolumn
     \bfseries\textit{\abstractname}:
   \else
     \begin{center}\textbf{\abstractname}\end{center}
   \fi}
  {\par}

\renewenvironment{IEEEkeywords}
  {\normalfont
   \@IEEEabskeysecsize
   \vspace{0.3em}
   \setlength{\parindent}{0pt}
   \noindent
   {\bfseries\itshape \IEEEkeywordsname: }\ignorespaces}
  {\par}

\def\IEEEbibitemsep{1.5pt plus .5pt}
\makeatother

\makeatletter
\AtBeginDocument{
  \renewcommand{\abstractname}{Tóm tắt}
}
\renewcommand{\IEEEkeywordsname}{Từ khóa}
\makeatother

\begin{document}
\pagenumbering{gobble}
\fontsize{9}{10}
\selectfont

\fancyhead[RE,LO]{\centering{\small{\textit{Hội thảo khoa học Quốc gia về Trí tuệ nhân tạo 2027 (FJCAI) - Cần Thơ, 27-28/3/2027}}}}

% ---------------------------------------------------------------------
% CHANGE 2 of 4: title and authors
% ---------------------------------------------------------------------
\title{Do Context Clusters Help Lightweight Image Classification?\\ An Ablation Study and a Stronger All-Local Baseline}

\author{\IEEEauthorblockN{Nguyen Viet Hung\IEEEauthorrefmark{1}, Nguyen Ba Nam\IEEEauthorrefmark{1}, Le Hoang Thai\IEEEauthorrefmark{1}}
	\IEEEauthorblockA{\IEEEauthorrefmark{1}Faculty of Information Technology, University of Science, VNU-HCM, Vietnam}
	\IEEEauthorblockA{Email: 23122032@student.hcmus.edu.vn, 23122043@student.hcmus.edu.vn, lhthai@fit.hcmus.edu.vn}
}

\maketitle

\begin{abstract}
% ~250 words. No special characters, symbols or math in title or abstract.
TODO: paste the final abstract here.
\end{abstract}

\begin{IEEEkeywords}
\textbf{\textit{lightweight image classification; ablation study; context cluster; knowledge distillation; efficient architectures}}
\end{IEEEkeywords}

\IEEEpeerreviewmaketitle

% ---------------------------------------------------------------------
% CHANGE 3 of 4: six-section body
% ---------------------------------------------------------------------
\section{Introduction}
\label{sec:intro}
% ~850 words, 7 paragraphs. Last paragraph = 2-sentence limitations preview.

\section{Related Work}
\label{sec:related}
% ~480 words, 3 paragraphs.
\subsection{Lightweight convolutional networks}
\subsection{Token mixers and MetaFormer}
\subsection{Knowledge distillation}

\section{Methodology}
\label{sec:method}
% ~750 words. Figure 1 goes at the top of this page.
\subsection{The hybrid block and its design space}
\subsection{The All-Local configuration}

\begin{figure*}[!t]
\centering
\includegraphics[width=\textwidth]{fig1_architecture.pdf}
\caption{The All-Local (2.7M) architecture. A stride-1 stem preserves the
$32\times32$ resolution into stage~1, whose token mixer is a learned
$3\times3$ convolution; stages~2--4 use $7\times7$ depthwise followed by
$1\times1$ pointwise convolution. Neither context clustering nor the fixed
LBP filters of the hybrid model appear in the final design.}
\label{fig:arch}
\end{figure*}

\section{Experiments}
\label{sec:exp}
% =====================================================================
% Section IV: Experiments.  \input{} from main.tex inside \section{Experiments}.
% Every number here is traceable to a run in configs/run_plan.yaml; cost figures
% come from the single-session M1-v3 benchmark.
% =====================================================================

\subsection{Experimental Setup}
\label{sec:setup}

\textbf{Datasets.} We evaluate on Tiny-ImageNet \cite{tinyimagenet} (200
classes) and CIFAR-10 (10 classes). Both are processed at $32\times32$
resolution. For Tiny-ImageNet this means downsampling from the native
$64\times64$, which discards information and makes our absolute numbers
\emph{not} comparable to work that evaluates at $64\times64$; we adopt it
deliberately, because the regime this paper studies is low-resolution
classification under a tight parameter budget, and we hold it fixed across
every model so that all comparisons reported here are internally valid.
Tiny-ImageNet does not release labels for its test split, so we use the
official 10{,}000-image validation split as our \emph{test} set and hold out
a further 10{,}000 images from the training split, drawn with a fixed seed,
as the validation set used for model selection. No model ever selects an
epoch using the reported test set.

\textbf{Budget parity.} Every model, proposed or baseline, is trained for the
same number of epochs on a given dataset: 100 on Tiny-ImageNet and 250 on
CIFAR-10. We use identical augmentation (RandAugment, random erasing, mixup
and CutMix), the same optimiser and schedule, and the same batch size. No
model receives a tuned recipe of its own.

\textbf{Knowledge distillation.} Where distillation is used, every student,
\emph{including every baseline}, is distilled from a ResNet-18 teacher
\cite{resnet} trained on the same dataset and the same seed, with the same
loss: $\alpha = 0.5$ and temperature $T = 4$ \cite{hinton}. We report both
cross-entropy (CE) and distilled (KD) results for all models, so that no
comparison mixes training signals.

\textbf{Seeds and statistics.} Headline configurations are trained with three
seeds ($42, 43, 44$); tables report the mean and standard deviation, with the
number of runs $n$ given explicitly. Differences between groups are tested
two ways. At the seed level we use Welch's unequal-variance $t$-test on the
three per-seed accuracies. Within a seed we use an exact paired McNemar test
on the $10{,}000$ test predictions of the two models being compared, which
uses the pairing that a $t$-test over seed means discards. We report both
because they answer different questions, and we state $n$ wherever it is
below three.

\textbf{Reproducibility.} Every run writes a provenance record containing the
git commit, the seed, the resolved teacher checkpoint and the full software
and hardware environment, together with its per-image test predictions. All
configurations, logs and predictions are released.

% ---------------------------------------------------------------------
\subsection{Main Results}
\label{sec:main}

Table~\ref{tab:tin} reports Tiny-ImageNet. Under distillation the proposed
All-Local model reaches $61.21 \pm 0.40$ top-1, against $59.72 \pm 0.14$ for
the parameter-matched ShuffleNetV2-x1.5 \cite{shufflenetv2} trained
identically from the same teacher. The margin of $+1.49$ pp is significant at
the seed level (Welch $t = 6.10$, $p = 0.016$) and, independently, on each of
the three seeds under a paired McNemar test. The same model exceeds the
ResNet-18 teacher that supervised it by $+2.46$ pp ($p = 0.008$) while using
$4.16\times$ fewer parameters and $2.41\times$ fewer multiply--accumulate
operations. Table~\ref{tab:c10} shows the comparison holds on CIFAR-10, where
the proposed model leads ShuffleNetV2-x1.5 by $+0.82$ pp with distillation
(McNemar $p = 0.009$) and by $+0.64$ pp without it ($p = 0.014$), and is
statistically indistinguishable from the teacher ($-0.08$ pp validation,
$+0.15$ pp test, $p = 0.432$) at $4.25\times$ fewer parameters.

\textbf{The cross-entropy picture is weaker, and we state it plainly.}
Without distillation the proposed model reaches $58.69 \pm 0.10$ against
$58.22 \pm 0.23$ for ShuffleNetV2-x1.5. That $+0.47$ pp margin does not reach
significance at the seed level ($p = 0.051$), and against the ResNet-18
teacher the proposed model is $-0.07$ pp, a tie ($p = 0.405$). The
beat-the-teacher result reported above therefore holds \emph{in the distilled
setting}, and we condition every such claim on that setting rather than
presenting it as an unconditional property of the architecture.

\textbf{An asymmetry we observe but do not explain.} Given the identical
distillation recipe, the proposed model gains $+2.52$ pp from distillation
while ShuffleNetV2-x1.5 gains $+1.50$ pp and ShuffleNetV2-x1.0 gains $+1.58$
pp. The proposed design converts the same teacher signal into roughly one
additional percentage point. We note that this is not explained by headroom:
the proposed model has the \emph{higher} cross-entropy accuracy of the pair
and therefore less room to improve, yet gains more. We have not tested a
mechanism for this and do not propose one; we report it as an observation and
leave it to future work. What the comparison does establish is that under a
fixed, identical distillation protocol applied to every model, the proposed
architecture is the more accurate one at a matched parameter budget.

\begin{table}[!t]
\renewcommand{\arraystretch}{1.15}
\caption{\uppercase{Tiny-ImageNet top-1 (\%), $32\times32$, 100 epochs for every model.
Mean $\pm$ standard deviation over $n$ seeds. Every KD student uses the same
ResNet-18 teacher, $\alpha = 0.5$, $T = 4$.}}
\label{tab:tin}
\centering
\footnotesize
\begin{tabular}{@{}lrrcc@{}}
\toprule
Model & Params & MACs & CE ($n$) & KD ($n$) \\
\midrule
ResNet-18 (teacher) & 11.27M & 556.8M & $58.76${\scriptsize$\pm$0.06} (3) & -- \\
MobileNetV2 \cite{mobilenetv2} & 2.48M & 25.8M & $49.25$ (1) & $52.32$ (1) \\
ShuffleNetV2-x1.0 & 1.46M & 46.3M & $55.78${\scriptsize$\pm$0.30} (3) & $57.36${\scriptsize$\pm$0.40} (3) \\
ShuffleNetV2-x1.5 & 2.68M & 96.1M & $58.22${\scriptsize$\pm$0.23} (3) & $59.72${\scriptsize$\pm$0.14} (3) \\
Context Cluster \cite{coc} & 2.46M & 34.4M & $49.45$ (1) & -- \\
Hybrid (ours, $\rho{=}0.5$) & 1.80M & 138.9M & $56.40${\scriptsize$\pm$0.21} (3) & $58.19${\scriptsize$\pm$0.25} (3) \\
\midrule
\textbf{All-Local (ours)} & 2.71M & 231.3M & $\mathbf{58.69}${\scriptsize$\pm$0.10} (3) & $\mathbf{61.21}${\scriptsize$\pm$0.40} (3) \\
\bottomrule
\end{tabular}
\end{table}

\begin{table}[!t]
\renewcommand{\arraystretch}{1.15}
\caption{\uppercase{CIFAR-10 top-1 (\%), $32\times32$, 250 epochs, single seed.}}
\label{tab:c10}
\centering
\footnotesize
\begin{tabular}{@{}lrcc@{}}
\toprule
Model & Params & CE & KD \\
\midrule
ResNet-18 (teacher) & 11.17M & $96.76$ & -- \\
MobileNetV2 & 2.24M & $94.06$ & $94.28$ \\
ShuffleNetV2-x1.0 & 1.26M & $95.40$ & $95.22$ \\
ShuffleNetV2-x1.5 & 2.49M & $95.56$ & $95.86$ \\
Hybrid (ours, $\rho{=}0.5$) & 1.76M & $95.44$ & $96.16$ \\
\midrule
\textbf{All-Local (ours)} & 2.65M & $\mathbf{96.20}$ & $\mathbf{96.68}$ \\
\bottomrule
\end{tabular}
\end{table}

% ---------------------------------------------------------------------
\subsection{Ablation Study}
\label{sec:ablation}

Table~\ref{tab:ablation} varies one design decision at a time from the hybrid
model, holding the stem, fold partitioning, block depth and width fixed so
that each row isolates a single factor.

\textbf{Clustering does not pay for itself.} Setting $\rho = 0$, which removes
the local convolutional branch and leaves pure clustering while keeping every
other choice identical, costs $-0.74$ pp. Setting $\rho = 1$ instead, removing
the clustering branch entirely, \emph{gains} $+1.03$ pp while using 10\% fewer
parameters. The two rows bracket the question the paper asks: at this scale and
resolution the local branch is what carries the accuracy, and the clustering
branch is a net cost.

\textbf{The supporting design choices contribute little.} Removing fold
partitioning ($+0.20$ pp) and removing channel shuffle ($+0.25$ pp) both leave
accuracy unchanged within run-to-run variation, so neither is load-bearing.
Changing cluster occupancy in either direction hurts ($-1.02$ pp at $r = 1$,
$-0.80$ pp at $r = 4$), and giving stage~4 four centres instead of one costs
$-0.31$ pp. The single largest effect in the table is the stem stride: moving
from stride~1 to stride~2 costs $-5.72$ pp. Because that change halves the
resolution entering every stage as well as altering cluster occupancy, we
attribute the bulk of it to resolution rather than to clustering.

\textbf{The hand-designed local operator is replaceable.}
Table~\ref{tab:operator} compares the fixed sparse LBP-pattern filters used at
stage~1 of the hybrid model against a plain learned $3\times3$ convolution at
matched parameters. At 1.62M parameters the learned convolution is ahead by
$0.10$ pp (McNemar $p = 0.63$); at the 2.71M scale of the proposed model the
two are indistinguishable over three seeds ($61.21 \pm 0.40$ against
$61.23 \pm 0.28$, Welch $p = 0.956$). The learned convolution nevertheless
uses $2.50\times$ less peak activation memory at batch 128, because the fixed
filters expand the channel count eightfold before projecting back. We
therefore adopt the learned convolution: the hand-designed operator buys no
accuracy and costs memory.

\begin{table}[!t]
\renewcommand{\arraystretch}{1.15}
\caption{\uppercase{One-change ablations on Tiny-ImageNet, cross-entropy.
Each row changes exactly one design decision from the hybrid baseline.}}
\label{tab:ablation}
\centering
\footnotesize
\begin{tabular}{@{}lccr@{}}
\toprule
Configuration & Top-1 (\%) & $n$ & $\Delta$ \\
\midrule
Hybrid baseline ($\rho{=}0.5$) & $56.40${\scriptsize$\pm$0.21} & 3 & -- \\
\midrule
$\rho = 0$ (clustering only) & $55.66${\scriptsize$\pm$0.20} & 3 & $-0.74$ \\
$\rho = 1$ (local only, 1.62M) & $57.43${\scriptsize$\pm$0.14} & 3 & $\mathbf{+1.03}$ \\
\midrule
No fold partitioning & $56.60$ & 1 & $+0.20$ \\
No channel shuffle & $56.65$ & 1 & $+0.25$ \\
Stage-4 $2\times2$ proposals & $56.09$ & 1 & $-0.31$ \\
Occupancy $r = 1$ & $55.38$ & 1 & $-1.02$ \\
Occupancy $r = 4$ & $55.60$ & 1 & $-0.80$ \\
Stem stride $1 \rightarrow 2$ & $50.68$ & 1 & $\mathbf{-5.72}$ \\
\bottomrule
\end{tabular}
\end{table}

\begin{table}[!t]
\renewcommand{\arraystretch}{1.15}
\caption{\uppercase{Stage-1 operator: fixed LBP-pattern filters against a learned
$3\times3$ convolution, at matched parameters. Memory is peak activation
memory at batch 128.}}
\label{tab:operator}
\centering
\footnotesize
\begin{tabular}{@{}llccr@{}}
\toprule
Scale & Stage-1 operator & Top-1 (\%) & $n$ & Mem. \\
\midrule
\multirow{2}{*}{1.62M, CE} & fixed LBP & $57.31$ & 1 & 796.9\,MB \\
 & learned $3\times3$ & $57.41$ & 1 & 316.9\,MB \\
\midrule
\multirow{2}{*}{2.71M, KD} & fixed LBP & $61.23${\scriptsize$\pm$0.28} & 3 & 1037.0\,MB \\
 & \textbf{learned} $\mathbf{3\times3}$ & $61.21${\scriptsize$\pm$0.40} & 3 & \textbf{415.0\,MB} \\
\bottomrule
\end{tabular}
\end{table}

% ---------------------------------------------------------------------
\subsection{Cost Analysis}
\label{sec:cost}

All cost figures in Table~\ref{tab:cost} are measured on a single NVIDIA
Tesla T4 (fp32, $32\times32$ input, 200-class head). Every model was profiled
in one session on the same host, so cross-model differences are not
confounded by hardware; we verified that re-profiling the same models in a
separate session leaves parameters, operation counts and peak memory
bit-identical but shifts batch-128 latency by up to 14\%, and we therefore
report a single session throughout. Latency is given both in eager execution
and with the forward pass captured in a CUDA Graph; every captured graph was
verified to reproduce its eager output bitwise, and each figure is the median
of five timed loops of 100 iterations after a warm-up of at least one second.
Peak memory is reset per batch size, and resident weights are reported
separately from activations.

At batch size one, the deployment-relevant operating point, the proposed model
occupies 10.37\,MB of weights and 22.58\,MB of peak memory, against 43.03\,MB
and 58.75\,MB for the ResNet-18 teacher: a $4.15\times$ reduction in model
size and a $2.60\times$ reduction in peak memory, together with $2.41\times$
fewer operations and a $2.18\times$ lower graph-captured latency (0.804\,ms
against 1.752\,ms), while being 2.46\,pp more accurate. Against the
parameter-matched ShuffleNetV2-x1.5 the proposed model uses $2.41\times$ more
operations, but the measured deployment cost is close: 22.58\,MB against
20.85\,MB of peak memory at batch one ($1.08\times$) and 0.804\,ms against
0.752\,ms of graph-captured latency ($1.07\times$). In eager execution the
ordering reverses and the proposed model is $1.72\times$ faster (3.50\,ms
against 6.03\,ms), because 87.5\% of ShuffleNetV2-x1.5's eager batch-one
latency is CPU kernel-launch overhead against 77.1\% for ours, a consequence
of its per-block channel split, shuffle and concatenation. At batch 128 the
proposed model remains the more expensive of the two, using 415.0\,MB against
133.0\,MB and reaching 3803 against 5768 images per second; we report this
because it is a regime our claims do not cover.

Two measurements caution against treating operation counts as a latency proxy.
First, realized arithmetic throughput at batch 128 varies by more than
fourfold across the models evaluated: ResNet-18 performs $2.41\times$ more
operations than the proposed model yet completes a batch of 128 in 28.06\,ms
against 32.76\,ms, because dense $3\times3$ convolutions map onto far
better-optimised kernels than the bandwidth-bound depthwise--pointwise stacks
used by lightweight designs. Second, the relative ordering of two models can
invert between execution modes, as the ShuffleNetV2-x1.5 comparison above
shows. We therefore report both modes and draw deployment conclusions only
where they agree. Fig.~\ref{fig:pareto} places every model on the accuracy
against parameters and accuracy against operations planes; the proposed model
is on the Pareto frontier of the first and is not Pareto-dominant on the
second, and we claim only the former.

\begin{table*}[!t]
\renewcommand{\arraystretch}{1.15}
\caption{\uppercase{Measured cost on one NVIDIA Tesla T4, fp32, $32\times32$ input,
200-class head, all models profiled in a single session. Memory is peak
allocated; latency is the median of five timed loops.}}
\label{tab:cost}
\centering
\footnotesize
\begin{tabular}{@{}lrrrrrrrr@{}}
\toprule
& & & \multicolumn{2}{c}{Peak memory (MB)} & \multicolumn{3}{c}{Latency (ms)} & Throughput \\
\cmidrule(lr){4-5}\cmidrule(lr){6-8}
Model & Params & MACs & $b{=}1$ & $b{=}128$ & eager $b{=}1$ & graph $b{=}1$ & graph $b{=}128$ & (img/s) \\
\midrule
ResNet-18 (teacher) & 11.27M & 556.8M & 58.75 & 326.7 & 2.64 & 1.752 & 28.06 & 4709 \\
MobileNetV2 & 2.48M & 25.8M & 19.62 & 124.3 & 4.74 & 0.674 & 12.85 & 9970 \\
ShuffleNetV2-x1.0 & 1.46M & 46.3M & 15.72 & 95.6 & 6.07 & 0.680 & 13.44 & 9683 \\
ShuffleNetV2-x1.5 & 2.68M & 96.1M & 20.85 & 133.0 & 6.03 & 0.752 & 22.03 & 5768 \\
Hybrid (ours, $\rho{=}0.5$) & 1.80M & 138.9M & 19.70 & 481.6 & 7.68 & 1.212 & 36.47 & 3408 \\
\midrule
\textbf{All-Local (ours)} & 2.71M & 231.3M & 22.58 & 415.0 & 3.50 & 0.804 & 32.76 & 3803 \\
\bottomrule
\end{tabular}
\end{table*}

\begin{figure*}[!t]
\centering
\includegraphics[width=\textwidth]{fig2_pareto.pdf}
\caption{Tiny-ImageNet accuracy against parameters (a) and against
multiply--accumulate operations (b), both on logarithmic axes. Accuracy is
the mean over three seeds; cost is measured on a Tesla T4 in a single
session. Every baseline shown is distilled from the same ResNet-18 teacher;
the teacher itself is trained with cross-entropy and is marked separately.
The proposed model lies on the parameter frontier but not on the operation
frontier.}
\label{fig:pareto}
\end{figure*}


\section{Discussion and Limitations}
\label{sec:limits}
% ~400 words. Cover: 32x32 only and incomparability to 64x64 work; the
% 100/250-epoch budget; 2.41x the MACs of ShuffleNetV2-x1.5; 3.12x its
% batch-128 memory and 0.66x its throughput; CIFAR-10 at a single seed;
% and the untested distillation asymmetry of Section IV-B.

\section{Conclusion}
\label{sec:conclusion}
% ~190 words, 4 sentences, no new claims.

\section*{Acknowledgement}
The authors would like to thank...

% ---------------------------------------------------------------------
% CHANGE 4 of 4: references. {99} reserves label width for 2-digit numbers.
% ---------------------------------------------------------------------
\balance
\begin{thebibliography}{99}

\bibitem{coc} X. Ma, Y. Zhou, H. Wang, C. Qin, B. Sun, C. Liu and Y. Fu,
``Image as set of points,'' in \emph{Proc. Int. Conf. Learning
Representations (ICLR)}, 2023.

\bibitem{metaformer} W. Yu, M. Luo, P. Zhou, C. Si, Y. Zhou, X. Wang,
J. Feng and S. Yan, ``MetaFormer is actually what you need for vision,''
in \emph{Proc. IEEE/CVF Conf. Computer Vision and Pattern Recognition
(CVPR)}, 2022, pp. 10819--10829.

\bibitem{hinton} G. Hinton, O. Vinyals and J. Dean, ``Distilling the
knowledge in a neural network,'' in \emph{NIPS Deep Learning and
Representation Learning Workshop}, 2015.

\bibitem{shufflenetv2} N. Ma, X. Zhang, H.-T. Zheng and J. Sun,
``ShuffleNet V2: Practical guidelines for efficient CNN architecture
design,'' in \emph{Proc. European Conf. Computer Vision (ECCV)}, 2018,
pp. 116--131.

\bibitem{mobilenetv2} M. Sandler, A. Howard, M. Zhu, A. Zhmoginov and
L.-C. Chen, ``MobileNetV2: Inverted residuals and linear bottlenecks,''
in \emph{Proc. CVPR}, 2018, pp. 4510--4520.

\bibitem{convnext} Z. Liu, H. Mao, C.-Y. Wu, C. Feichtenhofer, T. Darrell
and S. Xie, ``A ConvNet for the 2020s,'' in \emph{Proc. CVPR}, 2022,
pp. 11976--11986.

\bibitem{replknet} X. Ding, X. Zhang, J. Han and G. Ding, ``Scaling up your
kernels to 31x31: Revisiting large kernel design in CNNs,'' in
\emph{Proc. CVPR}, 2022, pp. 11963--11975.

\bibitem{resnet} K. He, X. Zhang, S. Ren and J. Sun, ``Deep residual
learning for image recognition,'' in \emph{Proc. CVPR}, 2016, pp. 770--778.

\bibitem{dkd} B. Zhao, Q. Cui, R. Song, Y. Qiu and J. Liang, ``Decoupled
knowledge distillation,'' in \emph{Proc. CVPR}, 2022, pp. 11953--11962.

\bibitem{tinyimagenet} Y. Le and X. Yang, ``Tiny ImageNet visual
recognition challenge,'' Stanford CS231N Technical Report, 2015.

\bibitem{mcnemar} T. G. Dietterich, ``Approximate statistical tests for
comparing supervised classification learning algorithms,''
\emph{Neural Computation}, vol. 10, no. 7, pp. 1895--1923, 1998.

% TODO: add the remaining references (target ~28 total).

\end{thebibliography}

\end{document}
